\scanpage{350}
Sous ces conditions, pour tout $U\in\operatorname{Ob}\mathcal X$,
\[
\deleted{\Gamma(U}\quad
\inserted{C^*(U)=}\operatorname{Hom}(U,C^*)=\Gamma(U,C_U^*)
\]
est \deleted{un} \inserted{$k(U)$-}complexe

$\uline{k(U)}:$ \deleted{$\Gamma(U,k_U)$} anneau commutatif unitaire
$=$ \unsure{définit des structures, ce qui implique que}
\deleted{$H^*\Gamma(U,C_U^*)$}
\inserted{$H^*=H^*(C^*(U))$}
est muni d'une structure évidente d'algèbre \inserted{graduée}
(associat., unit., anticomm\unsure{utative}).
Si $C^*$ est une \uline{résolution} de $k$, on dispose pour $U$ d'un hom
\[
(1)\qquad H^*(C^*(U))\longrightarrow H^*(U,k_U).
\]

\margin{\uline{en fait}}
\unsure{Il est évident que c'est un} homomorphisme d'anneaux.
\unsure{C'est l'objet} \illegible{2} \uline{\unsure{th.}} \unsure{d'anneaux comm.}

\uline{Ex. 1} Cas où $C^*$ \unsure{est une algèbre diff. usuelle}.

\uline{Ex. 2} $\Phi$ \unsure{est un} \illegible{2}
\unsure{anneau commutatif}, \unsure{avec une} $k$-\unsure{linéarité}.
\unsure{Supposons} $k\longrightarrow\Phi$, \unsure{libre} \illegible{3}
section $T$ de $\Phi$. \unsure{Supposons que} $\unsure{i\cdot k=T}$
\unsure{on le résout par le} complexe

\diagram{\begin{tikzcd}[column sep=small]
0\arrow[r] & C^0\arrow[r]\arrow[d,equal] & C^1\arrow[r]\arrow[d,equal] & \cdots\arrow[r] & C^n\arrow[r]\arrow[d,equal] & \unsure{C^{n+1}}\arrow[r]\arrow[d,equal] & \cdots\\
0\arrow[r] & \Phi\arrow[r,"T_\wedge"'] & \Lambda^2\Phi\arrow[r,"T_\wedge"'] & \cdots\arrow[r] & \Lambda^{n+1}\Phi\arrow[r,"T_\wedge"'] & \Lambda^{n+2}\Phi\arrow[r] & \cdots
\end{tikzcd}}

est une \uline{résolution} de $k$, \unsure{on a}
\[
Z^i\simeq\Lambda^i\Psi\qquad(i\geq0)
\]

(\unsure{en supposant que} \illegible{2} \unsure{est un} \illegible{2},
\unsure{ce qui correspond au} \illegible{2} \unsure{complexe}
\diagram{\begin{tikzcd}
\unsure{0}\arrow[r] & \Lambda^{\unsure{i-1}}\Psi\arrow[r] & \Lambda^{i}\Phi\arrow[r] & \Lambda^{\unsure{i}}\Psi\arrow[r] & 0
\end{tikzcd}})

\mspage{351}
